% =============================================================
%  Kapitel 8 – Nächste Schritte / Offene Punkte
% =============================================================
\section{Nächste Schritte und offene Punkte}
\label{sec:naechste_schritte}

Die folgende Liste fasst die wichtigsten offenen Punkte zusammen,
priorisiert nach Relevanz für die Bewertung.

\subsection{Priorisierte To-do-Liste}

\begin{enumerate}
  \item \textbf{[P1 – Detailed Design]}
    \textbf{Activity-Diagramme für private Methoden ausarbeiten.} \\
    Mindestens für \texttt{checkReleaseGuard()} und
    \texttt{triggerEmergencyStop()} ein Activity-Diagramm erstellen
    \emph{oder} schriftlich begründen, warum der Zustandsautomat
    (Abschnitt~\ref{sec:estop_statemachine}) als alleinige
    dynamische Modellierung ausreicht.
    \todo{Activity-Diagramme erstellen oder Begründung formulieren.}

  \item \textbf{[P1 – Detailed Design]}
    \textbf{Vollständige Methodenliste mit Vor-/Nachbedingungen.} \\
    Alle privaten Methoden der Klasse \texttt{EmergencyStop} mit
    Preconditions, Postconditions und ggf.\ Invarianten spezifizieren
    (Design by Contract).
    \todo{Tabelle mit Pre-/Postconditions für alle privaten Methoden vervollständigen.}

  \item \textbf{[P2 – Architektur]}
    \textbf{Kurze schriftliche Begründung der Sprachauswahl (C++) finalisieren.} \\
    Der Abschnitt~\ref{sec:sprachauswahl} enthält bereits die wesentlichen
    Argumente. Mit dem Team abstimmen und auf 3--5 Sätze verdichten.
    \todo{Sprachauswahl-Begründung mit Team abstimmen und finalisieren.}

  \item \textbf{[P2 – Architektur]}
    \textbf{CAN-Datenklasse ausdetaillieren.} \\
    Die Klasse für CAN-Pakete (Struktur: ID, Payload, Timestamp) in einem
    eigenen Klassendiagramm modellieren und in Abschnitt~\ref{sec:klassendiagramme}
    einbinden.
    \todo{CAN-Datenklasse modellieren und Diagramm einbinden.}

  \item \textbf{[P3 – Qualität]}
    \textbf{Latenzzahlen ergänzen.} \\
    Typische CAN-Zykluszeit und LAN-Latenz aus Systemtests eintragen
    (Abschnitt~\ref{sec:deployment}).
    \todo{Messwerte aus Systemtest nachtragen.}

  \item \textbf{[P3 – Optional]}
    \textbf{Logger-Interface vollständig spezifizieren.} \\
    Vollständige Methodenliste für \texttt{Logger} (inkl.\ Log-Level-Enum,
    Puffer-Strategie, Ausgabe-Kanal).
    \todo{Logger-Interface ausarbeiten (optional, aber gern gesehen laut Aufgabenstellung).}

  \item \textbf{[P3 – Optional]}
    \textbf{Interfaces (\texttt{EnvironmentControl}) vollständig ausarbeiten.} \\
    Nicht nur eine einzelne Funktion, sondern alle Methoden und Signale
    des Interfaces spezifizieren.
    \todo{EnvironmentControl-Interface vollständig modellieren.}
\end{enumerate}

\subsection{Bekannte Annahmen (zur Überprüfung)}

Die folgenden Annahmen wurden getroffen und müssen mit dem Team
bzw.\ der Aufgabenstellung validiert werden:

\begin{itemize}
  \item \assumption{Emergency Stop läuft auf dem Arduino Mega Slave
        (nahe DriveControlUnit). Erste Vermutung laut Mitschrieb; zu bestätigen.}
  \item \assumption{CAN-Zykluszeit $< 5\,\text{ms}$; zu messen.}
  \item \assumption{Distanzschwelle für Rot-Zone: $\leq 15\,\mathrm{cm}$.}
  \item \assumption{CAN-Datenklasse kapselt ID, Payload und Timestamp.}
  \item \assumption{Logger schreibt auf seriellen Ausgang des Slave-Arduinos;
        kein persistenter Speicher.}
\end{itemize}
